%% paginas/2_pre_textuais/abstract.tex
%% Elemento Pré-Textual Obrigatório (Abstract/Resumo em Inglês) -- Conforme ABNT NBR 6028
%% Regra: Tradução fiel do resumo vernáculo. Parágrafo único, espaçamento simples,
%% de 150 a 500 palavras. Keywords separadas por pontos.

\begin{center}
  \textbf{ABSTRACT}
\end{center}

\vspace{0.8cm}

\begin{spacing}{1.0}
  \noindent
  \MakeUppercase{\meuautor}. \textbf{\meutitulo}. \meuano. \pageref{LastPage} f. \tipotrabalho\ (\grauacademicopretendido) -- \minhainstituicao, \meulocal, \meuano.
\end{spacing}

\vspace{0.8cm}

\begin{spacing}{1.0}
  \noindent
  This work aims to present a highly modular and structured LaTeX template for the creation of academic documents in compliance with the most recent standards of the Brazilian Association of Technical Standards (ABNT). Through file modularization, dividing components into pre-textual, textual, and post-textual folders, and isolating style data into external files, it becomes much easier to maintain the project, develop academic work cooperatively, and comply strictly with formatting rules. The model addresses the precise configuration of margins, standard typography, line spacing, long citations, and list of sections. With this clean and expandable structure, the Overleaf user can focus exclusively on the content of their research, while the LaTeX engine handles the mechanical formatting. The practical compilation results demonstrate productivity gains and reduction of common errors in undergraduate final projects, master's theses, and doctoral dissertations.
\end{spacing}

\vspace{1cm}

\noindent
\textbf{Keywords:} LaTeX. ABNT. Academic Template. Modularity. Document Formatting.
